\begin{name}
	{HÌNH HỌC}
	{ĐỀ LUYỆN TẬP ONLINE KHỐI 8 - TUẦN \thechapter}
\end{name}
\setcounter{ex}{0}\setcounter{bt}{0}
\Opensolutionfile{ans}[ans/ans-2-TT-11-NguyenKhuyen-HCM-24-03-2019]
\begin{luuy}
	\noindent{\bf MỘT SỐ KIẾN THỨC CẦN NHỚ:}\\
	\textbf{Nhận xét:} Hai tam giác đồng dạng thì tỉ số diện tích bằng bình phương tỉ số đồng dạng.\\
	\textbf{ Ví dụ:} Ta có $\triangle ABC\sim \triangle MNP$.\\
	Suy ra $\dfrac{S_{\triangle ABC}}{S_{\triangle ABC}}=\left(\dfrac{AB}{MN}\right)^2=\left(\dfrac{AC}{MP}\right)^2=\left(\dfrac{BC}{NP}\right)^2$.
\end{luuy}


%%==========Câu 1
\begin{ex}
	Bóng của một cột điện trên mặt đất có độ dài $4{,}5$ m. Cùng thời điểm đó, một thanh sắt cao $2{,}1$ m cắm vuông góc với mặt đất có bóng dài $0{,}6$ m. Tính chiều cao của cột điện.
\loigiai{
	\immini{
		Biểu diễn cột điện bằng đoạn thẳng $AB$ và cây sắt là đoạn thẳng $A'B'$ như hình vẽ.\\
		Gọi $AC$ và $A'C'$ lần lượt là biểu diễn cho độ dài bóng của cột điện và thanh sắt ở dưới mặt đất.\\
		Ta có các độ dài được biểu diễn tương ứng trong hình vẽ.\\
		Xét $\triangle ABC$ và $\triangle A'B'C'$ ta có
		\begin{itemize}
			\item $\widehat{BAC}=\widehat{B'A'C'}=90^\circ$ (cột điện và thanh sắt vuông góc với mặt đất).
			\item $\widehat{ACB}=\widehat{A'C'B'}$ (hai góc đồng vị và hai tia sáng mặt trời song song)
		\end{itemize}
		Vậy $\triangle ABC\sim \triangle A'B'C'$ (góc - góc).\\
		Suy ra
		\allowdisplaybreaks
		\begin{eqnarray*}
			&&\dfrac{AB}{A'B'}=\dfrac{AC}{A'C'}\\
			&\Leftrightarrow&AB=\dfrac{A'B'\cdot AC}{A'C'}=\dfrac{2,1\cdot 4,5}{0,6}=15.75\text{ (m)}.
		\end{eqnarray*}
		
	}{
		\begin{tikzpicture}[]
			\def\dodai{2.5};
			\path
				(0,0) coordinate (A)
				($(A)+(90:4.5*\dodai)$) coordinate (B)
				($(A)+(0:1.29*\dodai)$) coordinate (C)
				(4,0) coordinate (A')
				($(A')+(90:2.1*\dodai)$) coordinate (B')
				($(A')+(0:0.6*\dodai)$) coordinate (C')
			;
			\draw
				(A)--(B)--(C)--cycle
				(A')--(B')--(C')--cycle
				(A)--(C')
			;
			\draw (A)--(C) node [midway,sloped,below]{$4,5$ m};
			\draw (A')--(B') node [midway,sloped,above]{$2,1$ m};
			\draw (A')--(C') node [midway,sloped,below]{$0,6$ m};
			\foreach \x/\g in {A/-90,A'/-90,B/90,B'/90,C/-45,C'/-45}
			{
				\draw [fill=black](\x) circle (1.5pt) +(\g:3mm) node {$\x$};
			}
		\end{tikzpicture}
	}
}
\end{ex}

%%==========Câu 2
\begin{ex}
	Cho tam giác $ABC$ đồng dạng với tam giác $A’B’C’$ theo tỉ số $k$. Biết diện tích tam giác $ABC$ bằng $24$ cm$^2$. 
	\begin{listEX}
		\item Chứng minh $\dfrac{S_{A B C}}{S_{A'B'C'}}=k^2$.
		\item Tính diện tích tam giác $A’B’C’$ với $k=\dfrac{2}{3}$.
	\end{listEX}
	\loigiai{
		\begin{center}
			\begin{tikzpicture}
			\def\dd{1}
			\def\goc{73}
			\path
			(0,0) coordinate (B)
			($(B)+(0:6*\dd)$) coordinate (C)
			($(B)+(\goc:4*\dd)$) coordinate (A)
			($(C)+(0:1cm)$) coordinate (B')
			($(B')+(0:3*\dd)$) coordinate (C')
			($(B')+(\goc:2*\dd)$) coordinate (A')
			($(B)!(A)!(C)$) coordinate (H)
			($(B')!(A')!(C')$) coordinate (H')
			;
			\draw
			(A)--(B)--(C)--cycle
			(A')--(B')--(C')--cycle
			(A)--(H)
			(A')--(H')
			;
			\foreach \x/\g in {A/90,A'/90,B/-135,B'/-135,C/-45,C'/-45,H/-90,H'/-90}
			{
				\draw [fill=black](\x) circle (1.5pt) +(\g:3mm) node {$\x$};
			}
			\end{tikzpicture}
		\end{center}
	\begin{listEX}
		\item Hạ đường cao $AH$ của $\triangle ABC$ và đường cao $A'H'$ của $\triangle A'B'C'$ như hình vẽ.
		ta có $\triangle ABC$ đồng dạng với $\triangle A'B'C'$ theo tỉ số $k$ nên ta có
		$$\dfrac{AB}{A'B'}=\dfrac{AC}{A'C'}=\dfrac{BC}{B'C'}=k.$$
		Xét $\triangle ABH$ và $\triangle A'B'H'$ có
		\begin{itemize}
			\item $\widehat{AHB}=\widehat{A'H'B'}=90^\circ$ ($AH$ và $A'H'$ lần lượt là đường cao $\triangle ABC$ và $\triangle A'B'C'$).
			\item $\widehat{ABH}=\widehat{A'B'H'}$ ($\triangle ABC\sim \triangle A'B'C'$).
		\end{itemize}
		Vậy $\triangle ABH\sim \triangle A'B'H'$ (góc - góc).\\
		Suy ra $\dfrac{AH}{A'H'}=\dfrac{AB}{A'B'}=k$.\\
		Vậy ta có 
		$$\dfrac{S_{ABC}}{S_{A'B'C'}}=\dfrac{\dfrac{1}{2}AH\cdot BC}{\dfrac{1}{2}A'H'\cdot B'C'}=\dfrac{AH\cdot BC}{A'H'\cdot B'C'}=\dfrac{AH}{A'H'}\cdot\dfrac{AB}{A'B'}=k\cdot k=k^2.$$
		\item Áp dụng câu a ta có
		\allowdisplaybreaks
		\begin{eqnarray*}
			&&\dfrac{S_{A B C}}{S_{A'B'C'}}=k^2\\
			&\Leftrightarrow& \dfrac{S_{A B C}}{S_{A'B'C'}}=\left(\dfrac{2}{3}\right)^2\\
			&\Leftrightarrow&\dfrac{S_{A B C}}{S_{A'B'C'}}=\dfrac{4}{9}\\
			&\Leftrightarrow&S_{A'B'C'}=\dfrac{S_{A B C}}{\dfrac{4}{9}}\\
			&\Leftrightarrow&S_{A'B'C'}=\dfrac{9\cdot S_{A B C}}{4}\\
			&\Leftrightarrow&S_{A'B'C'}=\dfrac{9\cdot 24}{4}\\
			&\Leftrightarrow&S_{A'B'C'}=54.
		\end{eqnarray*}
		Vậy diện tích của $\triangle A'B'C'$ là $54$ cm$^2$.
	\end{listEX}
	}
\end{ex}

%%==========Câu 3
\begin{ex}
	Cho hình chữ nhật $ABCD$ có $AB=12$ cm, $BC=9$ cm, $H$ là hình chiếu của $A$ lên $BD$.
	\begin{listEX}	
		\item Chứng minh: $\triangle AHB \sim \triangle BCD$.
		\item Tính diện tích tam giác $AHB$.
		\item Gọi $O$ là giao điểm của $AC$ và $BD$. Tính diện tích tam giác $AHO$. 
	\end{listEX}
\loigiai{
	\immini{
		\begin{listEX}	
			\item Xét $\triangle AHB$ và $\triangle BCD$ có
			\begin{itemize}
				\item $\widehat{AHB}=\widehat{BCD}=90^\circ$ ($H$ là hình chiếu của $A$ lên $BD$ và $\widehat{BCD}$ là một góc của hình chữ nhật).
				\item $\widehat{ABH}=\widehat{BDC}$ (cùng phụ $\widehat{DBC}$).
			\end{itemize}
			Vậy $\triangle AHB\sim \triangle BCD$ (góc - góc).
		\end{listEX}
	}{
		\begin{tikzpicture}
			\def\dd{1};
			\path
				(0,0) coordinate (D)
				($(D)+(0:4*\dd)$) coordinate (C)
				($(D)+(90:3*\dd)$) coordinate (A)
				($(C)+(90:3*\dd)$) coordinate (B)
				($(B)!(A)!(D)$) coordinate (H)
				(intersection of A--C and B--D) coordinate (O)
			;
			\draw
				(A)--(B)--(C)--(D)--cycle
				(B)--(D)
				(A)--(H)
				(A)--(C)
				(B)--(D)
			;
			\foreach \x/\g in {A/135,B/45,C/-45,D/-135,H/-45,O/90}
			{
				\draw [fill=black](\x) circle (1.5pt) +(\g:3mm) node {$\x$};
			}
		\end{tikzpicture}
	}
	\begin{listEX}
		\setcounter{enumi}{1}
		\item Ta có vì $ABCD$ là hình chữ nhật nên $AD=BC=9$ cm và $CD=AB=12$ cm.\\
		Xét $\triangle BCD$ vuông tại $C$. Theo định lý Py-ta-go ta có
		\allowdisplaybreaks
		\begin{eqnarray*}
			&&BD^2=BC^2+CD^2\\
			&\Leftrightarrow&BD^2=9^2+12^2\\
			&\Leftrightarrow&BC^2=15^2\\
			&\Rightarrow&BD=15.
		\end{eqnarray*}
		Vì $\triangle AHB\sim \triangle BCD$ nên ta có
		\allowdisplaybreaks
		\begin{eqnarray*}
			&&\dfrac{S_{AHB}}{S_{BCD}}=\left(\dfrac{AB}{BD}\right)^2\\
			&\Leftrightarrow& \dfrac{S_{AHB}}{S_{BCD}}=\left(\dfrac{4}{5}\right)^2\\
			&\Leftrightarrow&\dfrac{S_{AHB}}{S_{BCD}}=\dfrac{16}{25}\\
			&\Leftrightarrow&S_{AHB}=\dfrac{16}{25}\cdot S_{BCD}\\
			&\Leftrightarrow&S_{AHB}=\dfrac{16}{25}\cdot \dfrac{1}{2}\cdot BC\cdot CD\\
			&\Leftrightarrow&S_{AHB}=\dfrac{16}{25}\cdot \dfrac{1}{2}\cdot 9\cdot 12\\
			&\Leftrightarrow&S_{AHB}=\dfrac{864}{25}\\
			&\Leftrightarrow&S_{AHB}=34{,}56.
		\end{eqnarray*}
		\item 
		Vì $ABCD$ là hình chữ nhật nên ta có $O$ là trung điểm của cả $AC$ và $BD$.\\
		Xét $\triangle AOB$ có
		$$S_{AOB}=\dfrac{1}{2}AH\cdot OB=\dfrac{1}{2}AH\cdot \dfrac{1}{2}\cdot DB=\dfrac{1}{2}S_{ABD}=\dfrac{1}{2}\cdot \dfrac{1}{2}\cdot9\cdot12=27\text{ cm$^2$}.$$
		Vậy $S_{AHO}=S_{AHB}-S_{AOB}=34,56-27=7,56$ cm$^2$.
	\end{listEX}
}
\end{ex}

%%==========Câu 4
\begin{ex}
	Cho tam giác $ABC$ có trung tuyến $AM$. Gọi $I$ là trung điểm của $AM$, $BI$ cắt $AC$ tại $E$. Gọi $F$ là trung điểm của $BE$. 
	\begin{listEX}
		\item Chứng minh $\triangle BFM$ đồng dạng $\triangle BEC$; $\triangle IFM$ đồng dạng với $\triangle IEA$.
		\item Tính tỉ số $\dfrac{AE}{AC}$.
	\end{listEX}
	
\loigiai{
	\immini{
		\begin{enumerate}
			\item Xét $\triangle BFM$ và $\triangle BEC$ ta có
			\begin{itemize}
				\item Góc $\widehat{FBM}$ chung.
				\item $\dfrac{BF}{BE}=\dfrac{BM}{BC}=\dfrac{1}{2}$ ($M$ và $F$ lần lượt là trung điểm của $BC$ và $BE$).
			\end{itemize}
			Vậy $\triangle BFM\sim \triangle BEC$ (cạnh - góc - cạnh).
			\end{enumerate}
	}{
		\begin{tikzpicture}
			\path 
				(0,0) coordinate (B)
				(2,4) coordinate (A)
				(6,0) coordinate (C)
				($(B)!1/2!(C)$) coordinate (M)
				($(A)!1/2!(M)$) coordinate (I)
				(intersection of A--C and B--I) coordinate (E)
				($(B)!1/2!(E)$) coordinate (F)
			;
			\draw 
				(A)--(B)--(C)--cycle
				(A)--(M)
				(B)--(E)
				(A)--(F)
				(F)--(M)
			;
			\foreach \x/\g in {A/90,B/-135,C/-45,M/-90,E/45,I/-45,F/135}
			{
				\draw [fill=black](\x) circle (1.5pt) +(\g:3mm) node {$\x$};
			}
		\end{tikzpicture}
	}
	\begin{listEX}
		\item[~] Xét $\triangle BEC$ có $M$ và $F$ lần lượt là trung điểm của $BC$ và $BE$ nên ta có $MF$ là đường trung bình của $\triangle BEC$ nên $FM\parallel EC$ và $FM=\dfrac{1}{2}EC$.\\
		Xét $\triangle IFM$ và $\triangle IEA$ có 
		\begin{itemize}
			\item $\widehat{AIE}=\widehat{MIF}$ (đối đỉnh).
			\item $\widehat{IEA}=\widehat{IFM}$ (sole-trong và $FM\parallel AC$)
		\end{itemize}
		Vậy $\triangle IFM\sim \triangle IEA$ (góc - góc).
		\setcounter{enumi}{1}
		\item Vì $\triangle IFM\sim \triangle IEA$ nên ta có $\dfrac{FM}{AE}=\dfrac{IM}{IA}=1$ (vì $I$ là trung điểm của $AM$).\\
		Vậy $FM=AE$.\\
		Suy ra $\dfrac{AC}{AE}=\dfrac{AC}{FM}=\dfrac{AE+EC}{FM}=\dfrac{FM+EC}{FM}=1+\dfrac{EC}{MC}=1+2=3$. Vậy $\dfrac{AE}{AC}=\dfrac{1}{3}$.
	\end{listEX}
}
\end{ex}

%%==========Câu 5
%\begin{ex}
%	Cho tam giác $ABC$ có $3$ góc nhọn. Vẽ ba đường cao $AD$,  $BE$, $CF$ cắt nhau tại $H$.  
%	\begin{listEX}
%		\item Chứng minh $\triangle AEB \sim\triangle AFC$.
%		\item Chứng minh $\widehat{AFE}=\widehat{ACB}$.
%		\item Chứng minh $HB \cdot HE = HF \cdot HC$ và $\widehat{FEB}=\widehat{FCB}$.
%		\item Gọi $M$ là trung điểm của $BC$, đường thẳng đi qua $H$ vuông góc với $MH$ cắt $AB$, $AC$ lần lượt tại $P$, $Q$. Chứng minh $HP = HQ$.
%	\end{listEX}
%\loigiai{
%	\immini{
%		\begin{listEX}
%			\item Xét $\triangle AEB$ và $\triangle AFC$ có
%			\begin{itemize}
%				\item $\widehat{BAC}$ chung.
%				\item $\widehat{AEB}=\widehat{AFC}=90^\circ$ ($BE$ và $CF$ là đường cao).
%			\end{itemize}
%			Vậy $\triangle AEB\sim \triangle AFC$ (góc - góc ).
%		\end{listEX}
%	}{
%		\begin{tikzpicture}
%			\path 
%			(0,0) coordinate (B)
%			(2,4) coordinate (A)
%			(6,0) coordinate (C)
%			($(B)!(A)!(C)$) coordinate (D)
%			($(A)!(C)!(B)$) coordinate (F)
%			($(A)!(B)!(C)$) coordinate (E)
%			(intersection of A--D and B--E) coordinate (H)
%			($(B)!1/2!(C)$) coordinate (M)
%			;
%			\draw 
%			(A)--(B)--(C)--cycle
%			(A)--(D)
%			(B)--(E)
%			(C)--(F)
%			(E)--(F)
%			(M)--(H)
%			;
%			\foreach \x/\g in {A/90,B/-135,C/-45,D/-90,E/45,F/135,H/-45}
%			{
%				\draw [fill=black](\x) circle (1.5pt) +(\g:3mm) node {$\x$};
%			}
%		\end{tikzpicture}
%	}
%	\begin{listEX}
%		\setcounter{enumi}{1}
%		\item Vì $\triangle AEB\sim \triangle AFC$ nên $\dfrac{AE}{AF}=\dfrac{AB}{AC}\Leftrightarrow\dfrac{AE}{AB}=\dfrac{AF}{AC}$.\\
%		Xét $\triangle AEF$ và $\triangle ABC$ có
%		\begin{itemize}
%			\item $\dfrac{AE}{AB}=\dfrac{AF}{AC}$ (chứng minh trên).
%			\item $\widehat{BAC}$ là góc chung.
%		\end{itemize}
%		Vậy $\triangle AEF\sim \triangle ABC$ (cạnh - góc - cạnh).\\
%		Vậy $\widehat{AFE}=\widehat{ACB}$ (hai góc tương ứng).
%		\item Xét $\triangle HFB$ và $\triangle HEC$ có
%		\begin{itemize}
%			\item $\widehat{FHB}=\widehat{EHC}$ (đối đỉnh).
%			\item $\widehat{HFB}=\widehat{HEC}=90^\circ$ ($CF$ và $BE$ lần lượt là đường cao).
%		\end{itemize}
%		Vậy $\triangle HFB\sim \triangle HEC$ (góc - góc).\\
%		Suy ra $\dfrac{HF}{HE}=\dfrac{HB}{HC}\Leftrightarrow HB\cdot HE=\cdot HF\cdot HC$.\\
%		Xét $\triangle HFE$ và $\triangle HCB$ ta có
%		\begin{itemize}
%			\item $\widehat{EHF}=\widehat{CHB}$ (đối đỉnh).
%			\item $\dfrac{HE}{HC}=\dfrac{HF}{HB}$ $\left(\dfrac{HF}{HE}=\dfrac{HB}{HC}\right)$.
%		\end{itemize}
%		Vậy $\triangle HFE\sim \triangle HCB$ (cạnh - góc - cạnh).\\
%		Suy ra $\widehat{FEB}=\widehat{FCB}$ (hai góc tương ứng).
%		\item Gọi $M$ là trung điểm của $BC$, đường thẳng đi qua $H$ vuông góc với $MH$ cắt $AB$, $AC$ lần lượt tại $P$, $Q$. Chứng minh $HP = HQ$.
%	\end{listEX}
%}
%\end{ex}
\begin{ex}
%	\immini{
		Cho tam giác $ABC$ có $3$ góc nhọn. Vẽ ba đường cao $AD$, $BE$, $CF$ cắt nhau tại $H$
		\begin{enumerate}
			\item Chứng minh $\triangle AEB\sim\triangle AFC$.
			\item Chứng minh $\widehat{AFE}=\widehat{ACB}$.
			\item Chứng minh $HB\cdot HE=HF\cdot HC$ và $\widehat{FEB}=\widehat{FCB}$.
			\item Gọi $M$ là trung điểm của $BC$, đường thẳng đi qua $H$ vuông góc với $MH$ cắt $AB$, $AC$ lần lượt tại $P$, $Q$. Chứng minh $HP=HQ$.
		\end{enumerate}
%	}{
%		\begin{tikzpicture}[scale=1, line join=round, line cap=round, >=stealth]
%		%Phan Khai Bao Mac Dinh
%		\tkzInit[ymin=-0.5,ymax=4.5,xmin=-0.5,xmax=6.5]
%		%\tkzAxeXY
%		%\tkzGrid
%		\tkzClip
%		\tkzSetUpPoint[color=black,fill=black,size=6]
%		%Dinh Hinh
%		\tkzDefPoints{0/0/B,2/4/A,6/0/C}
%		\tkzDrawPolygon(A,B,C)
%		\tkzDrawAltitude(B,C)(A)\tkzGetPoint{D}
%		\tkzDrawAltitude(A,C)(B)\tkzGetPoint{E}
%		\tkzDrawAltitude(A,B)(C)\tkzGetPoint{F}
%		\tkzDefMidPoint(B,C)\tkzGetPoint{M}
%		\tkzInterLL(A,D)(C,F)\tkzGetPoint{H}
%		\tkzDefLine[orthogonal=through H](H,M) \tkzGetPoint{p}
%		\tkzInterLL(H,p)(A,B)\tkzGetPoint{P}
%		\tkzInterLL(H,p)(A,C)\tkzGetPoint{Q}
%		\tkzMarkRightAngles[fill=gray!20](B,F,C A,D,C B,E,C)
%		\tkzDrawPoints(A,B,C,D,E,F,M,H,P,Q)
%		\tkzDrawSegments(E,F H,M P,Q)
%		\tkzLabelPoints[above](A)
%		\tkzLabelPoints[below](D,M)
%		\tkzLabelPoints[right](H)
%		\tkzLabelPoints[left]()
%		\tkzLabelPoints[above right](E,Q)
%		\tkzLabelPoints[above left](F)
%		\tkzLabelPoints[below right](C)
%		\tkzLabelPoints[below left](B,P)
%		\end{tikzpicture}
%	}
	\loigiai{
		\immini{
			\begin{enumerate}
				\item Xét $\triangle AEB$ vuông tại $E$ và $\triangle AFC$ vuông tại $F$ có
				\begin{itemize}
					\item $\widehat{AEB}=\widehat{AFC}=90^\circ$.
					\item Góc $\widehat{BAC}$ chung.
				\end{itemize}
				Vậy $\triangle AEB\sim\triangle AFC$ (góc - góc).
			\end{enumerate}
		}{
			\begin{tikzpicture}[scale=1, line join=round, line cap=round, >=stealth]
			%Phan Khai Bao Mac Dinh
			\tkzInit[ymin=-0.5,ymax=4.5,xmin=-0.5,xmax=6.5]
			%\tkzAxeXY
			%\tkzGrid
			\tkzClip
			\tkzSetUpPoint[color=black,fill=black,size=6]
			%Dinh Hinh
			\tkzDefPoints{0/0/B,2/4/A,6/0/C}
			\tkzDrawPolygon(A,B,C)
			\tkzDrawAltitude(B,C)(A)\tkzGetPoint{D}
			\tkzDrawAltitude(A,C)(B)\tkzGetPoint{E}
			\tkzDrawAltitude(A,B)(C)\tkzGetPoint{F}
			\tkzDefMidPoint(B,C)\tkzGetPoint{M}
			\tkzInterLL(A,D)(C,F)\tkzGetPoint{H}
			\tkzDefLine[orthogonal=through H](H,M) \tkzGetPoint{p}
			\tkzInterLL(H,p)(A,B)\tkzGetPoint{P}
			\tkzInterLL(H,p)(A,C)\tkzGetPoint{Q}
			\tkzMarkRightAngles[fill=gray!20](B,F,C A,D,C B,E,C)
			\tkzDrawPoints(A,B,C,D,E,F,M,H,P,Q)
			\tkzDrawSegments(E,F H,M P,Q)
			\tkzLabelPoints[above](A)
			\tkzLabelPoints[below](D,M)
			\tkzLabelPoints[right](H)
			\tkzLabelPoints[left]()
			\tkzLabelPoints[above right](E,Q)
			\tkzLabelPoints[above left](F)
			\tkzLabelPoints[below right](C)
			\tkzLabelPoints[below left](B,P)
			\end{tikzpicture}
		}
		\begin{enumerate}
			\setcounter{enumi}{1}
			\item Ta có $\triangle AEB\sim\triangle AFC \text{  (cmt)}\Rightarrow\dfrac{AE}{AF}=\dfrac{AB}{AC}\Rightarrow\dfrac{AE}{AB}=\dfrac{AF}{AC}$.\\
			Xét $\triangle AEF$ và $\triangle ABC$ có
			\begin{itemize}
				\item Góc $\widehat{BAC}$ chung.
				\item $\dfrac{AE}{AB}=\dfrac{AF}{AC}$. (chứng minh trên)
			\end{itemize}
			Vậy $\triangle AEF\sim\triangle ABC$ (cạnh - góc - cạnh).\\
			Vậy $\widehat{AFE}=\widehat{ACB}$ (hai góc tương ứng).
			\item Xét $\triangle HFB$ vuông tại $F$ và $\triangle HEC$ vuông tại $E$ có
			\begin{itemize}
				\item $\widehat{HFB}=\widehat{HEC}=90^\circ$.
				\item $\widehat{FHB}=\widehat{EHC}$. (đối đỉnh)
			\end{itemize}
			Vậy $\triangle HFB\sim\triangle HEC$ (góc - góc).\\
			Vậy ta có $\dfrac{HF}{HE}=\dfrac{HB}{HC}\Rightarrow HF\cdot HC=HE\cdot HB$.\\
			Ta có $\dfrac{HF}{HE}=\dfrac{HB}{HC}\Rightarrow \dfrac{HF}{HB}=\dfrac{HE}{HC}$.\\
			Xét $\triangle HEF$ và $\triangle HCB$ có
			\begin{itemize}
				\item $\widehat{EHF}=\widehat{CHB}$. (đối đỉnh)
				\item $\dfrac{HF}{HB}=\dfrac{HE}{HC}$. (chứng minh trên)
			\end{itemize}
			Vậy $\triangle EHF\sim\triangle CHB$ (cạnh - góc - cạnh).\\
			Vậy $\widehat{FEH}=\widehat{CH}$ (hai góc tương ứng) $\Rightarrow \widehat{FEB}=\widehat{FCB}$.
			\item 
			\textbf{\textit{Cách 1:}} 
			Ta có
			\begin{itemize}
				\item $\widehat{AHE}=\widehat{BHD}$ (đối đỉnh).
				\item $\widehat{AHE}+\widehat{EAH}=90^\circ$ (hai góc nhọn trong $\triangle AHE$ vuông tại $E$).
				\item $\widehat{BHD}+\widehat{HBD}=90^\circ$ (hai góc nhọn trong $\triangle BHD$ vuông tại $D$).
			\end{itemize}
			Vậy $\widehat{EAH}=\widehat{HBD}$.\\
			Ta có
			\begin{itemize}
				\item $\widehat{EHQ}=\widehat{PHB}$ (đối đỉnh).
				\item $\widehat{EHQ}+\widehat{EQH}=90^\circ$ (hai góc nhọn trong $\triangle EQH$ vuông tại $E$).
				\item $\widehat{PHB}+\widehat{BHM}=90^\circ$ ($PH\perp HM$).
			\end{itemize}
			Vậy $\widehat{BHM}=\widehat{AQH}$.\\
			Xét $\triangle HBM$ và $\triangle QAH$ có
			\begin{itemize}
				\item $\widehat{QAH}=\widehat{HBM}$.
				\item $\widehat{BHM}=\widehat{AQH}$.
			\end{itemize}
			Vậy $\triangle HBM\sim\triangle QAH$ (góc - góc). $\Rightarrow \dfrac{MH}{HQ}=\dfrac{MB}{AH}$.\hfill$(1)$\\
			Ta có
			\begin{itemize}
				\item $\widehat{AHF}=\widehat{CHD}$ (đối đỉnh).
				\item $\widehat{AHF}+\widehat{FAH}=90^\circ$ (hai góc nhọn trong $\triangle AHF$ vuông tại $F$).
				\item $\widehat{CHD}+\widehat{HCD}=90^\circ$ (hai góc nhọn trong $\triangle DHC$ vuông tại $D$).
			\end{itemize}
			Vậy $\widehat{FAH}=\widehat{HCD}$.\\
			Ta có
			\begin{itemize}
				\item $\widehat{FHP}=\widehat{CHQ}$ (đối đỉnh).
				\item $\widehat{FHP}+\widehat{FPH}=90^\circ$ (hai góc nhọn trong $\triangle FPH$ vuông tại $E$).
				\item $\widehat{CHQ}+\widehat{CHM}=90^\circ$ ($QH\perp HM$).
			\end{itemize}
			Vậy $\widehat{APH}=\widehat{CHM}$.\\
			Xét $\triangle HMC$ và $\triangle PHA$ có
			\begin{itemize}
				\item $\widehat{PAH}=\widehat{HMC}$.
				\item $\widehat{APH}=\widehat{CMH}$.
			\end{itemize}
			Vậy $\triangle HMC\sim\triangle PHA$ (góc - góc). $\Rightarrow \dfrac{MH}{HP}=\dfrac{MC}{AH}$.\hfill$(2)$\\
			Từ $(1)$ và $(2)$ ta có $\dfrac{MH}{HQ}=\dfrac{MB}{AH}=\dfrac{MC}{AH}=\dfrac{MH}{HP}\Rightarrow\dfrac{MH}{HQ}\dfrac{MH}{HP}\Rightarrow HP=HQ$.
			
			\immini{
				\textbf{\textit{Cách 2:}} Gọi $K$ là điểm đối xứng của $C$ qua $H$,\\
				Xét $\triangle CBK$ có
				\begin{itemize}
					\item $M$ là trung điểm của $BC$.
					\item $H$ là trung điểm của $CK$.
				\end{itemize}
			}{
				\begin{tikzpicture}[scale=1, line join=round, line cap=round, >=stealth]
				%Phan Khai Bao Mac Dinh
				\tkzInit[ymin=-0.5,ymax=4.5,xmin=-2.5,xmax=6.5]
				%\tkzAxeXY
				%\tkzGrid
				\tkzClip
				\tkzSetUpPoint[color=black,fill=black,size=6]
				%Dinh Hinh
				\tkzDefPoints{0/0/B,2/4/A,6/0/C}
				\tkzDrawPolygon(A,B,C)
				\tkzDrawAltitude(B,C)(A)\tkzGetPoint{D}
				\tkzDrawAltitude(A,C)(B)\tkzGetPoint{E}
				\tkzDrawAltitude(A,B)(C)\tkzGetPoint{F}
				\tkzDefMidPoint(B,C)\tkzGetPoint{M}
				\tkzInterLL(A,D)(C,F)\tkzGetPoint{H}
				\tkzDefLine[orthogonal=through H](H,M) \tkzGetPoint{p}
				\tkzInterLL(H,p)(A,B)\tkzGetPoint{P}
				\tkzInterLL(H,p)(A,C)\tkzGetPoint{Q}
				\tkzDefPointBy[rotation= center H angle 180](C) \tkzGetPoint{K}
				\tkzInterLL(K,P)(H,B)\tkzGetPoint{I}
				\tkzMarkRightAngles[fill=gray!20](B,F,C A,D,C B,E,C K,I,H)
				\tkzDrawPoints(A,B,C,D,E,F,M,H,P,Q,K,I)
				\tkzDrawSegments(E,F H,M P,Q K,H K,B K,I)
				\tkzLabelPoints[above](A,K)
				\tkzLabelPoints[below](D,M)
				\tkzLabelPoints[right](H)
				\tkzLabelPoints[left]()
				\tkzLabelPoints[above right](E,Q)
				\tkzLabelPoints[above left](F)
				\tkzLabelPoints[below right](C,I)
				\tkzLabelPoints[below left](B,P)
				\end{tikzpicture}
			}
			\noindent Vậy $MH$ là đường trung bình trong tam giác $CBK$. $\Rightarrow HM\parallel KB$.\\
			Ta có $\heva{&HM\parallel KB\\&HM\perp PQ}\Rightarrow PQ\perp KB$.\\
			Xét $\triangle KHB$ có
			\begin{itemize}
				\item $HP$ là đường cao. ($PQ\perp KB$)
				\item $BF$ là đường cao. ($BF\perp KH$)
				\item $HP$ cắt $BF$ tại $P$
			\end{itemize}
			Vậy $P$ là trực tâm của $\triangle KBH$. $\Rightarrow KP\perp BH$. Gọi giao điểm  của $KP$ và $HB$ là $I$.\\
			Xét $\triangle HIK$ vuông tại $I$ và $\triangle HEC$ vuông tại $E$ có
			\begin{itemize}
				\item $\widehat{IHK}=\widehat{EHC}$ (đối đỉnh).
				\item $HK=HC$ (do cánh vẽ điểm $K$)
			\end{itemize}
			Vậy $\triangle HIK=\triangle HEC$ (cạnh huyền - góc nhọn).$\Rightarrow HI=HE$ (hai cạnh tương ứng).\\
			Xét $\triangle HIP$ vuông tại $I$ và $\triangle HEQ$ vuông tại $E$ có
			\begin{itemize}
				\item $\widehat{IHP}=\widehat{EHQ}$ (đối đỉnh).
				\item $HI=HE$ (chứng minh trên)
			\end{itemize}
			Vậy $\triangle HIP=\triangle HEQ$ (cạnh góc vuông - góc nhọn).$\Rightarrow HP=HQ$ (hai cạnh tương ứng).
			
		\end{enumerate}
	}
\end{ex}
%%==========Câu 6
\begin{ex}
	Cho tam giác $ABC$, đường cao $AH$. Biết $AH=\sqrt{HB\cdot HC}$. Chứng minh $\triangle ABC$ vuông tại $A$.
\loigiai{
	\noindent{\it Cách 1.}
	\immini{
		Ta có
			\begin{eqnarray*}
				&&AH=\sqrt{HB\cdot HC}\\
				&\Leftrightarrow&AH^2=HB\cdot HC\\
				&\Leftrightarrow&\dfrac{AH}{HC}=\dfrac{HB}{AH}.
			\end{eqnarray*}
		Xét $\triangle AHB$ và $\triangle CHA$ có
		\begin{itemize}
			\item $\widehat{AHB}=\widehat{AHC}=90^\circ$ ($AH$ là đường cao).
			\item $\dfrac{AH}{HC}=\dfrac{HB}{AH}$ (chứng minh trên).
		\end{itemize}
	}{
		\begin{tikzpicture}
		\path 
		(0,0) coordinate (B)
		(2,4) coordinate (A)
		(6,0) coordinate (C)
		($(B)!(A)!(C)$) coordinate (H)
		;
		\draw
		(A)--(B)--(C)--cycle
		(A)--(H)
		;
		\foreach \x/\g in {A/90,B/-135,H/-90,C/-90}
		{
			\draw [fill=black](\x) circle (1.5pt) +(\g:3mm) node {$\x$};
		}
		\end{tikzpicture}
	}
Vậy $\triangle AHB\sim \triangle CHA$ (cạnh - góc - cạnh).\\
Suy ra $\widehat{HAB}=\widehat{HCA}$.\\
Mà 
\begin{eqnarray*}
	&&\widehat{HAB}+\widehat{HBA}=90^\circ\\
	&\Rightarrow&\widehat{HCA}+\widehat{HBA}=90^\circ\\
	&\Rightarrow&\widehat{BAC}=90^\circ\\
\end{eqnarray*}
Vậy $\triangle ABC$ vuông tại $A$.\\
	\noindent{\it Cách 2.}\\
	\immini{
		Qua $A$ kẻ đường vuông góc với $AB$ cắt $BC$ ở $C'$.\\
		Xét $\triangle AHB$ và $\triangle C'HA$ có
		\begin{itemize}
			\item $\widehat{AHB}=\widehat{AHC'}=90^\circ$ ($AH$ là đường cao).
			\item $\widehat{BAH}=\widehat{AC'H}$ (cùng phụ $\widehat{HAC}$).
		\end{itemize}
	}{
		\begin{tikzpicture}
			\path 
				(0,0) coordinate (B)
				(2,4) coordinate (A)
				(6,0) coordinate (C)
				($(B)!(A)!(C)$) coordinate (H)
				($(A)!1!90:(B)$) coordinate (c)
				(intersection of A--c and B--C) coordinate (C')
			;
			\draw
				(A)--(B)--(C)--cycle
				(A)--(H)
				(A)--(C')
				(B)--(C')
			;
						\foreach \x/\g in {A/90,B/-135,C'/-45,H/-90,C/-90}
						{
							\draw [fill=black](\x) circle (1.5pt) +(\g:3mm) node {$\x$};
						}
		\end{tikzpicture}
	}
	\noindent Vậy $\triangle AHB\sim \triangle C'HA$ (góc - góc).\\
	Suy ra $\dfrac{AH}{HC'}=\dfrac{HB}{AH}\Leftrightarrow AH^2=HB\cdot HC'\Leftrightarrow AH=\sqrt{HB\cdot HC'}$.
	Vậy lúc này $C$ phải là điểm trùng với $C'$ hay $\triangle ABC$ vuông tại $A$.
}
\end{ex}
\centering{---Hết---}
\Closesolutionfile{ans}